\uuid{QPhY}
\titre{Sous-espace vectoriel/Représentation cartésienne}
\auteur{hanani ; mzari}
\organisation{exo7}
\licence{CC BY-NC-SA 4.0}
\datecreate{2020-01-20}
\niveau{L1}
\module{Algèbre}
\chapitre{Espace vectoriel}
\sousChapitre{Base}
\difficulte{3}
\theme{Base et dimension, Sous-espace vectoriel, Représentation cartésienne}
\exo7qcmid{lille-2.espaces-vectoriels.042}

\contenu{
\begin{question}
\qclassification{106.04}
\qtags{motcle=Sous-espace vectoriel/Représentation cartésienne}
Dans l'espace vectoriel $\Rr^3$, on pose $u_1=(1,0,1), u_2=(-1,1,1)$, $u_3=(1,-1,0)$ et on considère les sous-espaces vectoriels $E=\mbox {Vect}\{u_1,u_2\}$ et $F=\mbox {Vect}\{u_3\}$. Quelles sont les assertions vraies ?
\begin{answers}
\good{$E$ est un plan vectoriel}.
\bad{Une équation cartésienne de $E$ est $x+2y+z=0$}.
\good{$F$ est une droite vectorielle}.
\bad{Une équation cartésienne de $F$ est $z=0$}.
\end{answers}
\begin{explanations} $E$ est un plan vectoriel.
Soit $M(x,y,z)$ un vecteur de $\Rr^3$. $M \in E$ si et seulement s'il existe $a,b \in \Rr$ tels que $M=au_1+bu_2$. En résolvant ce système, on obtient une équation cartésienne de $E$ : $x+2y-z=0$.
\vskip2mm
$F$ est une droite vectorielle ; c'est donc l'intersection de deux plans de $\Rr^3$. Soit $M(x,y,z)$ un vecteur de $\Rr^3$. $M \in F$ si et seulement s'il existe un réel $a$ tels que $M=au_3$. En
résolvant ce système, on obtient une représentation cartésienne de $F$ : $(\mathtt{S})
\left\{\begin{array}{rcc}x+y&=&0\\
z&=&0.\end{array}\right.$
\end{explanations}
\end{question}
}
